\scanpage{88}
% manuscript lines 1–4
il s'ensuit que $\mathcal C$ est \uline{cohérent}. D'ailleurs, $\mathcal C$ satisfait $(S_2)$ en tout point de $Z$, car \deleted{si $x\in\unsure{Z}$, $V$ ouvert $\subset$ on a $\unsure{V|\mathcal M}\subset X$, on a $\Gamma(V,\mathcal C)\longrightarrow\Gamma(V\cap U,\mathcal C)$ est un isom.} par construction $\mathcal C\overset{\sim}{\longrightarrow}i_*\!i^*\mathcal C$.

% manuscript lines 4–9
Comme $\mathcal C\simeq B$ aux pts de $U$, \deleted{il s'} et que $B_x$ est une algèbre finie et normale sur $\underline O_{X,x}$, donc $(S_2)$ comme \unsure{anneau}, et par suite comme module, il s'ensuit que $\mathcal C$ est $(S_2)$ partout \inserted{comme Module,}. Celà \deleted{signifie} \unsure{signifie} aussi que l'anneau $C$ qui définit $\mathcal C$ est $S_2$. Bien \unsure{entendu} $A\subset C\subset K$.

% manuscript lines 5–9
% slanted
\margin{\unsure{Inutile} le fait que, $A$ étant quotient d'un régulier, on a $\dim A_{\mathfrak q}=\dim B_{\mathfrak p}$ si $\mathfrak p\in\operatorname{Spec} B$ induit $\mathfrak q$] cf verso}

% manuscript lines 9–18
Reste à prouver que $C$ est \deleted{\unsure{un|normal}} $(R_1)$. Soit $\mathfrak p$ un idéal premier de $C$ tel que $\dim C_{\mathfrak p}=1$, prouvons $C_{\mathfrak p}$ régulier. Soit $\mathfrak q$ l'idéal induit par $\mathfrak p$ sur $A$, alors [$A$ \deleted{est} quotient d'un régulier] on a $\dim A_{\mathfrak q}=1$, donc ou bien $A_{\mathfrak q}$ est normal, donc $A_{\mathfrak q}=B_{\mathfrak q}=B_{\mathfrak p}$ \deleted{\unsure{est donc}} normal, ou bien $A_{\mathfrak q}$ non normal, i.e. un des $\mathfrak p_i$, mais alors on a encore par construction $B_{\mathfrak q}$ normal, donc \deleted{$\unsure{B_{\mathfrak q}}$} $B_{\mathfrak p}$ est normal comme localisé de $B_{\mathfrak q}$. Celà achève la démonstration

% manuscript lines 19–23
\deleted{$A$ anneau (loclt noeth \unsure{par ex}), $B$ alg. finie sur $A$, \inserted{tt les idéaux \illegible{2}} $M$ module de type fini sur $B$, je dis que $M$ est $(S_n)$ sur $A$ ssi il est $(S_n)$ sur $B$. En \unsure{fait} supposons $A\subset B$, $\operatorname{Supp}M=\operatorname{Sp}B$, donc $f(X)=Y$. Soit $y\in Y$, on a $\operatorname{prof} M_y=\operatorname{Inf}_{z\,\mathrm{sur}\,y}\operatorname{prof}M_z$ \unsure{$\Rightarrow$} si $M$ sur $B$ est $(S_n)$, on a $\operatorname{prof}M_y\geq\operatorname{Inf}$}

% manuscript lines 22–23
\diagram{\begin{tikzcd}[ampersand replacement=\&,row sep=small]
f\&X\arrow[d]\\ \&Y
\end{tikzcd}}
% archive sheet number: 87

